\chapter{Watson--Crick and biochemical automata}
Two readers move along a double strand. One waits while the other advances.
Can that delay store information even when the controller has only three
states? Yes: the distance between the readers is part of the configuration.
This observation separates a two-head formal automaton from both an ordinary
finite-state reader and a biochemical device that consumes its input.

Chapter 16 developed languages and molecular recognition. Here we ask what an
automaton's state actually contains. We build a recognizer, prove what it
recognizes, and then examine a molecular mechanism with a different contract.
The Python results are calculations, not measurements of DNA reactions.

\section{Three meanings of state}
\Term{Control state} is a label such as \texttt{seek}. A complete computational
configuration also includes reader positions and the input. A molecular state
includes chemical species, concentrations and exposed binding sites. Confusing
these meanings can turn a modest finite-state experiment into an unsupported
claim about unbounded computation.

Our input is a pair of equal-length words $(u,v)$ whose symbols satisfy a
positionwise relation $\rho$. For a fixed input, a configuration is
\begin{equation}
 c=(q,i,j),\qquad 0\leq i,j\leq n,
\end{equation}
where $i$ and $j$ count consumed symbols on the upper and lower words.
A rule $(q,x,y,q')$ consumes upper block $x$ and lower block $y$ if both
match at their respective cursors. Either block may be empty, but not both.
The rule changes the control state to $q'$ and advances by $(|x|,|y|)$.

\Plate{01-heads}{A duplex has chemical polarity; an abstract automaton has
coordinate conventions. Both formal readers move to the right in this chapter,
although the illustrated DNA backbones are antiparallel. This is a reading
model, not a depiction of two polymerases synthesizing in opposite directions.}{fig:heads}

\Term{Watson--Crick automata} use paired inputs and independently moving
readers. We adopt the positionwise pairing and full-input acceptance convention
described by Czeizler and colleagues \cite{czeizler}, with two explicit choices:
we include the empty word and exclude rules consuming nothing on both strands.
Our examples use an abstract identity relation on $\{a,b\}$, not literal
adenine--adenine base pairing. A chemical encoding would need separate design
and validation. Formal complementarity need not mean the canonical DNA alphabet.

\section{Use head lag to recognize equal blocks}
We want the language $L=\{a^k b^k:k\geq0\}$. A one-head finite automaton
cannot remember an arbitrarily large $k$. Two cursors can compare the two
blocks without a growing list of control states. Set $u=v$, start at
$(\texttt{seek},0,0)$, and use these rules:
\begin{center}
\begin{tabular}{llll}\toprule
From & Upper & Lower & To\\\midrule
seek & $\varepsilon$ & a & seek\\
seek & a & b & pair\\
pair & a & b & pair\\
pair & b & $\varepsilon$ & drain\\
drain & b & $\varepsilon$ & drain\\\bottomrule
\end{tabular}\end{center}
Accepting control states are \texttt{seek} and \texttt{drain}, but acceptance
also requires $i=j=n$. The start state therefore accepts only the empty input
without transitions. Reaching an accepting label early is not sufficient.

\Plate{02-lag}{The accepting path for aabb has three phases. Lower-only
motion first creates a lag, paired motion compares the blocks, and upper-only
motion drains the remaining suffix. Coordinates, not the names of the three
control states, retain the input-dependent information.}{fig:lag}

\subsection{Prove both directions before writing a benchmark}
For $a^k b^k$ with $k>0$, lower-only steps consume the first $k$ lower
symbols. Exactly $k$ paired steps then consume upper $a$ symbols and lower
$b$ symbols. Finally $k$ upper-only steps consume the upper $b$ suffix.
Both cursors finish, so every desired input has an accepting run.

Conversely, a nonempty accepting run must enter \texttt{pair} and then
\texttt{drain}. The upper word contains only the $a$ symbols consumed in
\texttt{pair}, followed by the $b$ symbols consumed in \texttt{drain}.
The lower word contains only the $a$ symbols consumed in \texttt{seek},
followed by the $b$ symbols consumed in \texttt{pair}. Because the words
are identical, their $a$ counts agree and their $b$ counts agree. The paired
steps consume one of each, forcing all three phase lengths to equal $k$.
Thus acceptance implies membership in $L$. No appeal to a molecular speedup
was needed for this proof.

\section{Execute a configuration search}
The reference performs breadth-first exploration. Its queue contains only
configurations already reached by valid transitions. The parent map serves
both as a visited set and as a certificate from which one accepting path can
be reconstructed. The full companion file supplies \texttt{RULES}, \texttt{RHO}
and imports; the listing is extracted directly from that tested file.
\Listing{paired_run}

For aabb, the generated certificate is
\begin{center}\input{\Generated/trace.tex}\end{center}
Every successive row must correspond to one declared rule; the final row must
consume both words. Merely checking the returned Boolean would miss a broken
certificate or an acceptance test that ignores a stranded cursor.

For arbitrary block rules, matching a transition also costs work proportional
to the compared block lengths in this simple implementation. Counting visited
configurations alone does not account for that cost, queue storage, parent
records or input preparation. A simulator's work is also not the elapsed time
of a hypothetical molecular realization of the same transition relation.

\section{Nondeterminism does not manufacture free witnesses}
When $\rho$ is not injective, one upper word may admit several lower words.
The existential language convention accepts $u$ if \emph{some} admissible
$v$ has an accepting run. The following helper enumerates those candidates;
it does not promise an efficient way to find a useful one.
\Listing{lower_words}

\Plate{03-witness}{An existential lower strand is a witness, not an oracle.
With two admissible partners at each of four positions, 16 paired inputs are
possible. A formal accepting path does not specify how to synthesize, select
or reliably detect the successful member of a molecular pool.}{fig:witness}

If position $r$ has $d_r$ possible partners, the number of admissible lower
words is $\prod_{r=1}^{n}d_r$. When each $d_r=2$, this becomes $2^n$.
Parallel exploration changes the organization of the search, not that count.
An accepting species can be a small fraction of a mixture even when its
logical existence is enough for the language definition. Readout must
distinguish that species from background and unintended products.

Our tiny test uses $\rho$ containing both $a\leftrightarrow x$ and
$a\leftrightarrow y$, with a rule that accepts only the pair $(a,y)$.
The lower word $x$ is admissible but not successful. This separates pairing
validity, transition applicability and final acceptance into three tests.

\section{What the biochemical finite automaton implements}
Benenson and colleagues implemented a different, biochemical finite automaton
using DNA, a restriction nuclease and ligase \cite{benenson}. An exposed sticky
end encodes a current state and symbol. A compatible transition molecule binds
and is ligated; restriction cleavage exposes the next encoded end. The
transition molecule's geometry determines the cleavage frame. These are
chemical encodings of a transition, not two independently addressable heads.

\Plate{04-cleavage}{Conceptual restriction-based transition cycle. Recognition
and ligation create a substrate on which cleavage exposes the next encoded
end and releases a fragment. The recognition site and cleavage position are
distinct. Lengths, base sequences and enzyme-specific offsets are intentionally
not specified; this plate is not a laboratory protocol.}{fig:cleavage}

This mechanism gives a useful implementation checklist. Which exposed end
represents $(q,a)$? Which transition molecules compete for it? What molecule
represents $(q',\text{next symbol})$? Which products are waste rather than
future inputs? How is a terminal end converted into an observable reporter?
A correct answer to the abstract transition table answers none of these
chemical questions automatically.

The reported device also separates computation from detection: terminal
products interact with output detectors whose reporters can be distinguished.
Consequently, a missing signal can mean no accepting computation, incomplete
chemistry, insufficient reporter formation or inadequate detection. A chemical
failure and a rejected word must not silently share one Boolean variable.
The formal recognizer in this chapter neither reproduces that experiment nor
predicts its yields. Its role is to make the computational semantics explicit.

\section{Bound the search without hiding the heads}
For fixed words of length $n$, the generic configuration bound is
\begin{equation}
 N_{\rm config}\leq |Q|(n+1)^2.
\end{equation}
Each rule increases $i+j$ by at least one, so every run has at most $2n$
transitions. This monotone rank rules out cycles even if the control graph
contains self-loops. The reference still uses a visited set to merge paths
that reach the same configuration.

\Plate{05-frontier}{Generated exploration counts for $a^k b^k$ lie below
the generic configuration bound. This compares counts in a Python model,
not reaction time, molecular energy or physical memory consumption.}{fig:frontier}

Each cursor requires enough distinguishable values to represent $0$ through
$n$: at least $\lceil\log_2(n+1)\rceil$ bits in a binary address model.
Calling the controller finite-state is true; calling the complete configuration
constant-size is false. The read-only input also occupies $2n$ abstract symbols.
For our successful example, the path has $3k$ transitions, while exploration
may additionally visit dead ends. The upper bound and the actual path length
answer different questions.

\section{Validation and failure analysis}
The tests enumerate every binary word through length nine and compare the
recognizer with an independent direct language predicate. They also validate
every returned transition, reverse rule order, reject incompatible pairings,
and check the configuration bound. Invalid double-empty transitions are
rejected at the interface. These tests support the finite implementation;
the language argument above supplies the all-length reasoning.

A useful extension is to allow a double-empty rule. The monotone-rank proof
then fails immediately, although a finite visited-set search can still
terminate for fixed input. Another extension lets readers move left. Now the
$2n$ path bound fails too. Such changes are new models, not harmless options.

\section{Exercises with worked answers}
\begin{enumerate}
\item Trace ab. \textbf{Answer:} $(s,0,0),(s,0,1),(p,1,2),(d,2,2)$,
where $s,p,d$ abbreviate the three states. Both cursors end after three rules.
\item Why reject aab? \textbf{Answer:} Only one lower $b$ is available for
pairing. One upper $a$ remains, which \texttt{drain} cannot consume.
\item Why reject abb? \textbf{Answer:} After one pair step an unmatched
lower $b$ remains. Upper-only draining cannot finish the lower cursor.
\item Is the empty word accepted accidentally? \textbf{Answer:} No; the
initial configuration is explicitly final and fully consumed. Remove
\texttt{seek} from the final set to change the language to $k>0$.
\item Bound configurations at $n=10$. \textbf{Answer:} $3(11)^2=363$,
including unreachable combinations. This is not a claim that all are visited.
\item Bound any path at $n=10$. \textbf{Answer:} At most 20 transitions,
because total consumed length begins at zero and cannot exceed 20.
\item Count a successful path for $k=10$. \textbf{Answer:} Thirty rules
on words of length 20; this is below the general bound of 40.
\item How many witnesses with partner counts $(2,1,3,2)$?
\textbf{Answer:} Twelve. Summing the counts would incorrectly treat independent
position choices as mutually exclusive whole-word alternatives.
\item Does reversing rule order change the language? \textbf{Answer:} Not
for exhaustive search. It may change which certificate is returned first.
\item What does a sticky end encode here? \textbf{Answer:} In the described
biochemical architecture, a state--symbol interface. It is not a freely
addressable cursor into both strands.
\item Can absence of fluorescence prove rejection? \textbf{Answer:} Not
without a calibrated reporter, positive controls and a detection model.
Several chemical and measurement failures can produce the same observation.
\item Why is this not a constant-memory solution to equal-block recognition?
\textbf{Answer:} The state labels are constant in number, but cursor positions
and the supplied input scale with $n$. The model has changed the resources
available to the recognizer.
\end{enumerate}

\section{Vocabulary and the next question}
\Term{Configuration}: all information needed to determine legal next steps.
\Term{Head lag}: difference between the two consumed-prefix lengths.
\Term{Witness}: an admissible auxiliary input whose existence establishes
acceptance under an existential convention.
\Term{Cleavage frame}: position of cleavage relative to an encoded symbol,
used in the biochemical architecture to distinguish subsequent states.

The next chapter asks a stricter question: what additional storage and
simulation guarantees justify a claim of universality? A family of small
automata, an unbounded tape model and a finite molecular experiment must not
be merged into one claim.
